\section{Introduction}
\label{sec:intro}

The rapid proliferation of AI security frameworks---from NIST AI RMF
and MITRE ATLAS to OWASP LLM Top~10, the EU AI Act, and ISO/IEC
42001---has created a compliance landscape in which organizations must
simultaneously satisfy multiple overlapping standards. Each framework
uses its own taxonomy, control identifiers, and terminology, making it
difficult for practitioners to determine which controls in one framework
correspond to controls in another.

Manual crosswalk construction is the current industry practice. Subject
matter experts review pairs of controls across frameworks and assign
equivalence judgments. This process is slow (a single 12-framework
crosswalk can take hundreds of hours), subjective (inter-annotator
agreement is often low for partial-overlap cases), and brittle (each
framework revision invalidates portions of the mapping).

We propose an automated crosswalk system that scales to arbitrary
framework pairs. Our pipeline ingests framework control texts, builds a
multi-framework knowledge graph, and trains a supervised classifier to
assign one of four relationship tiers to each candidate control pair.
The system incorporates three signal families: (1)~lexical similarity
via BM25, (2)~dense semantic similarity via BGE-large embeddings, and
(3)~structural graph features extracted from a GATv2 attention network
trained on the full control graph.

A Mondrian conformal prediction wrapper provides finite-sample coverage
guarantees per class, and a KL-divergence router identifies predictions
where the model is uncertain, routing them to human reviewers. The
result is a human-in-the-loop system that automates the bulk of
crosswalk construction while preserving expert oversight where it
matters most.

We evaluate on 400 expert-labeled frozen test pairs across 26 framework
pairs, achieving tier accuracy of 0.37 with a bootstrap 95\% confidence
interval of [0.33, 0.42]. On LLM-labeled validation data the system
reaches 0.65 accuracy, highlighting a significant distribution gap
between expert and LLM labeling that we analyze in detail.
